% Rascunho de complemento à seção de trabalhos correlatos do TCC II.
% Preservar os trabalhos já citados no TCC I e integrar estas referências
% após receber o .tex e o .bib originais. As chaves abaixo estão em
% referencias-selecionadas-clack.bib.

Um sistema de controle de acesso a laboratório foi desenvolvido com ESP32 e
identificação por radiofrequência (RFID), acompanhado de uma aplicação web de
gerenciamento. O dispositivo se comunica com o servidor pelo protocolo HTTP,
e os dados são armazenados em um banco MySQL. No trabalho, a aplicação e o
banco foram hospedados no serviço Amazon EC2 \cite{goncalves2019rfidweb}.

Outra proposta reúne ESP32, etiquetas RFID e um servidor Apache, com uma
aplicação em PHP e banco de dados MySQL. O sistema também utiliza um aplicativo
móvel para o cadastro e a remoção de etiquetas \cite{costa2020esp32rfid}.

Em contexto de instituição de ensino, o IF Access foi proposto para controlar
o acesso aos ambientes do Instituto Federal da Paraíba por meio de RFID e
microcontrolador. O trabalho relata simulações envolvendo o gerenciamento de
usuários, ambientes e horários de acesso \cite{silva2022ifaccess}.

Assim como os trabalhos apresentados, o Clack integra identificação por RFID
ao gerenciamento informatizado de acessos. Na versão desenvolvida, o servidor
é executado localmente e a tranca utiliza permissões armazenadas para continuar
operando durante a indisponibilidade de rede, conforme as condições de
sincronização do sistema. A verificação de horários de aula permanece como
possibilidade de desenvolvimento futuro.
